\chapter{系统实现}
\label{chap:system}

本章介绍将第\ref{chap:method}章方法落地为可训练系统的工程实现，重点是\textbf{无侵入注入}、\textbf{推理与采样链路}、以及\textbf{奖励裁判服务}。

\section{无侵入的持续学习注入}

本文以 verl（0.8.0）作为底层强化学习框架，并遵循一条硬性工程原则：\textbf{不 fork、不改动框架源码}。第\ref{chap:method}章的复合损失（式\ref{eq:cl-loss}）与缓冲区钩子均通过框架提供的\textbf{外部注入接口}挂载：

\begin{itemize}
  \item \textbf{损失注入}：在框架初始化 worker 之后，调用 \texttt{set\_loss\_fn} 接口把复合损失闭包注入 actor。闭包不捕获缓冲区（缓冲区仅存于驱动进程，不随闭包序列化到 worker）。
  \item \textbf{缓冲区钩子}：包裹框架的参数更新入口，在每个训练步的\textbf{更新前}把携带梯度的回放行拼接进 actor 批次，\textbf{更新后}回收本步产生的优胜轨迹入桶、并回填抗遗忘优先级。
  \item \textbf{外部模块注入}：少量必须在每个 worker（含推理副本进程）生效的补丁，通过框架的外部模块加载机制（环境变量）在 import 时挂载，同样不触碰框架源码。
\end{itemize}

这种设计使本文的持续学习逻辑与框架版本解耦：框架升级时只需合入官方提交，本文的注入层无需改动。

\section{训练指标落盘}

除框架原生指标外，本文额外落盘框架未提供的\textbf{离散度指标}以监控训练健康度：每步的奖励标准差、优势标准差、以及 GRPO \textbf{组内}奖励标准差（组内轨迹奖励全同意味着无学习信号）。这些指标经同一指标通道写入 JSONL，用于事后分析策略是否坍缩、优势估计是否退化。

\section{推理与采样链路}

\textbf{推理引擎}采用 LightLLM 提供高并发的策略采样。\textbf{基座模型}设计上面向更大规模，当前实现与验证以 Qwen3.5-9B（混合注意力架构）进行。

\textbf{工具执行}在沙箱中完成。沙箱后端通过统一的客户端契约与注册表解耦——同一套调用代码可切换本地（开发）与云端（生产）后端，换厂商零改调用方。

\textbf{纯文本约束}：本文训练聚焦纯文本能力域。当工具产出图片等多模态内容时，含图请求会在网关层被拦截丢弃，避免打崩纯文本推理路径。

\section{奖励裁判服务}

奖励裁判（第\ref{sec:reward}节）作为\textbf{外部服务}调用，而非框架内置的奖励模型 worker。裁判模型不写死在代码里，由配置/环境变量解析（裁判服务地址与模型名），可本地部署也可接远程端点。裁判开启思考模式，需为其推理预算预留足够的输出 token，避免中途截断导致 JSON 解析失败。对候选裁判模型，本文用统一细则、多次重复打分对比其打分\textbf{稳定度}与响应延迟，据此选型。
